

\begin{figure}
\noindent\begin{minipage}{.48\textwidth}
\begin{lstlisting}[language=Verus,style=VerusLineNos]
struct InvCell<T> { /* ... */ }

impl InvCell<T> {
    // Well-formedness of the InvCell
    #[spec] pub fn wf(&self) -> bool;

    // Boolean predicate indicating the values allowed to be stored.
    #[spec] pub fn inv(&self, val: T) -> bool;

    // Construct a new InvCell, with initial value `val` and invariant given by `f`.
    #[exec] pub fn new(val: T, #[spec] f: impl Fn(T) -> bool) -> Self {
        requires(f(val));
        ensures(|cell: Self| cell.wf() && forall(|t: T| f(t) == cell.inv(t)));
        /* ... */
    }
\end{lstlisting}
\end{minipage}
\hfill\begin{minipage}{.48\textwidth}
\begin{lstlisting}[language=Verus,style=VerusLineNos,firstnumber=17]
    // Write to the cell and return the old value.
    #[exec] pub fn replace(&self, val: T) -> T {
        requires(self.wf() && self.inv(val)); %*\clnum\label{clnum:invcell:replace:pre}*
        ensures(|old_val: T| self.inv(old_val));
        /* ... */
    }

    // Read the current value of the cell.
    #[exec] pub fn get(&self) -> T
        where T: Copy
    {
        requires(self.wf());
        ensures(|val: T| self.inv(val));  %*\clnum\label{clnum:invcell:get:post}*
        /* ... */
    }
}
\end{lstlisting}
\end{minipage}
    \caption{
        \label{fig:inv-cell-api}
        API and specification for \texttt{InvCell<T>}.
    }
\end{figure}

\section{Supporting Interior Mutability} \label{sec:interior-mutability}

Interior mutability is a Rust pattern in which the contents (the ``interior'') of a datatype \verb`X` may be modified even when it is shared via a reference type \verb`&X`.
Since \verb`&` is supposed to be an ``immutable'' reference, interior mutability appears to be at odds with the core tenets of Rust's type system, and in fact
interior mutability is only sound when restricted appropriately.
Rust's standard library provides a handful of types with interior mutability, e.g., \verb`Cell`, \verb`RefCell`, \verb`RwLock`, each of which provides a different set of restrictions
and characteristics. For example, \verb`Cell` may not be shared across threads, while \verb`RwLock` is thread-safe but incurs all the costs of being a lock.
The most flexible Rust datatype supporting interior mutability is the \verb`UnsafeCell`, upon which the aforementioned types are implemented. %on top in Rust's standard library.
Since \verb`UnsafeCell` has no restrictions, it is---as the name implies---\emph{not} safe in general, and implementations that use it must take great care.

Providing safe and correct versions of such types in \verus is challenging, 
since in our SMT encoding, values of type \verb`&T` are always treated as immutable.
Therefore, to handle any \verb`Cell`-like datatypes, our SMT representation of \verb`&Cell<T>` cannot include an encoding of its mutable ``interior'' \verb`T`.
How, then, are we able to verify programs that require reasoning about this interior?

There are two broad classes of strategies a \verus developer can use:
\begin{enumerate}
  \item Use linear ghost state to represent the contents of a cell, similar to the way linear ghost state represents the value pointed to by a pointer.
  \item Avoid ``keeping track of'' the interior value entirely. Instead, when the interior value is read, model the result as being effectively nondeterministic,
      potentially using invariants to restrict the set of values that can be stored in the interior.
\end{enumerate}
\verus provides primitives and additional verified libraries supporting both styles, which the user can mix-and-match as needed.

The first strategy is the one used by our primitive \verb`PCell` (``permissioned \verb`Cell`'').
In the same way that \verb`PPtr` is our safe alternative to Rust's raw pointers,
\verb`PCell` is our safe alternative to \verb`UnsafeCell`.
\verb`PCell` uses a ghost permission mechanism with a similar API and specification to \verb`PPtr`, allowing us to track the interior value on the permission object.

The second strategy is exemplified by the type \verb`InvCell` of Verus' standard
library (\autoref{fig:inv-cell-api}), which provides a \verb`Cell`-like interface and allows the user to specify an invariant as a boolean predicate on values.
Whenever they write to the cell, they must prove the written value satisfies the invariant~\clref{clnum:invcell:replace:pre},
and when they read from it, they obtain an arbitrary value that they can \emph{assume} satisfies it~\clref{clnum:invcell:get:post}.
We first illustrate how this can be useful in our next section, and then we
discuss how \verb`InvCell` is itself verified in terms of lower-level invariant primitives.


\subsection{Verified Example: Memoized Function Calls} \label{sec:memoize}

\begin{figure}

\noindent\begin{minipage}{.5\textwidth}

\begin{lstlisting}[language=Verus,style=VerusLineNos]
#[spec] fn expected_result() -> u64 { /* ... */ }

#[exec] fn computation() -> u64 {
    ensures(|res: u64| res == expected_result()); %*\lclnum\label{clnum:comp:post}*
    /* ... */
}

#[spec] fn cell_value_inv(v: Option<u64>) -> bool {
    equal(v, Option::Some(expected_result()))
    || equal(v, Option::None) %*\lclnum\label{clnum:memo:cell-value-inv}*
}

#[spec] fn cell_is_valid(
        cell: InvCell<Option<u64>>) -> bool {
    cell.wf()
    && forall(|v| (#[trigger] cell.inv(v) ==
        cell_value_inv(v)))
}

#[exec] fn init_cell() -> InvCell<Option<u64>> {
    ensures(|c| cell_is_valid(c));
    InvCell::new(Option::None,
        |v: Option<u64>| cell_value_inv(v)) %*\lclnum\label{clnum:memo:init}*
}
\end{lstlisting}
\end{minipage}\hfill
\noindent\begin{minipage}{.5\textwidth}

\begin{lstlisting}[language=Verus,style=VerusLineNos]
#[exec] fn memoized_computation(
        cell: &InvCell<Option<u64>>) -> u64 {
    requires(cell_is_valid(*cell));
    ensures(|res: u64| res == expected_result());

    match cell.get() {
        Option::Some(res) => res, %*\lclnum\label{clnum:memo:case-none}*
        Option::None => {
            let res = computation();
            cell.replace(Option::Some(res));
            res %*\lclnum\label{clnum:memo:case-some}*
        }
    }
}

struct Client<'a> {
    cell: &'a InvCell<Option<u64>>,
}

fn main() {
    let c = init_cell();
    let client1 = Client { cell: &c };
    let client2 = Client { cell: &c };
    let x = memoized_computation(&client1.cell);
    let y = memoized_computation(&client2.cell);
    assert(x == y);
}
\end{lstlisting}
\end{minipage}

    \caption{
        \label{fig:memoize}
        Memoization example built on top of \texttt{InvCell}. %We take as given some (spec) value \texttt{expected\_result()} and
        %some (exec) code \texttt{computation()} that returns \texttt{expected\_result()}.
        %We then produce a function \texttt{memoized\_computation(cell)} which also returns
        %\texttt{expected\_result()}.
    }
\end{figure}

\emph{Memoization} is an optimization technique whereby a user saves time
by storing the result of a computation the first time it is invoked; on future invocations,
they use the stored value. Here, we show how to memoize a function call
\texttt{computation()}.
In \autoref{fig:memoize} we use a function that takes 0 arguments for simplicity, so there is only a single value to memoize;
in our supplementary materials~\cite{verussupplementary}, we provide a slightly more complex example that memoizes
a single-argument function \texttt{computation(i)}.

To set up the problem, we assume that \verb`computation()` has a postcondition~\lclref{clnum:comp:post} ensuring
that its result is equal to some desired (spec-mode) value, \verb`expected_result()`.
(This is similar to the setup of \verb`fibo_impl` and \verb`fibo` from earlier.)
The aim is to construct a function \verb`memoized_computation` that also returns
\verb`expected_result()`.
To keep the problem interesting, we also insist that it be possible to share the 
``result store'' across potentially many clients.
As such, we need to use a shared reference type \verb`&`; however,
a given update invocation might need to update the result store, which requires
mutability. Therefore, we need to use some form of interior mutability.

In our approach, we use
an \verb`InvCell` with a simple invariant on the data held by the cell.
When initializing the cell~\lclref{clnum:memo:init}, we specify the data invariant as a (spec-mode) boolean predicate on the interior values;
here, we set it to the function \verb`cell_value_inv`, defined at \lclref{clnum:memo:cell-value-inv}.
The resulting property of the cell is expressed as in \verb`cell_is_valid`.
This definition says that the value is valid if and only if the stored value is either
\verb`None` (not yet computed) or contains the correct answer.

To implement \verb`memoized_computation`, we first read from the cell; if the value we get is
\verb`Some`, then we return the value immediately~\lclref{clnum:memo:case-none} (as we can assume it satisfies the invariant
we just specified).
Otherwise, we perform the computation, store it in the cell,
and return it~\lclref{clnum:memo:case-some}.

Finally, in \verb`main`, we show that we can create multiple ``clients,'' sharing a
reference to the cell, and use them to call \verb`memoized_computation`.



\subsection{Invariant Primitives and \texttt{InvCell} Verification} \label{sec:invcell}

Just as Rust's standard library implements \verb`Cell` via \verb`UnsafeCell`,
in \verus we can implement and verify \verb`InvCell` via our 
\verb`UnsafeCell`-equivalent, \verb`PCell`.
To do this, though, we first need to introduce our \emph{invariant primitives}.

To see what these are for, consider what happens when we
try to implement \verb`InvCell<T>` using \verb`UnsafeCell<T>`.
From the API, we know that we need to be able to write
even when we only have access to a shared reference \verb`&InvCell<T>`,
but writing to the underlying \verb`UnsafeCell<T>` requires exclusive ownership
of the \verb`PermData<T>` object.

Once again, we run into this problem of trying to gain exclusive ownership of something
that is shared. However, we have pushed the problem one layer down---to the \emph{ghost}
layer, and this is where \verus introduces its trusted invariant primitives to escape
the problem.

The two primitives are called \verb`LocalInvariant<G>` and
\verb`AtomicInvariant<G>`.\footnote{Both these types also have additional type parameters used to specify the invariant as a boolean predicate on \texttt{G}.}
Each one allows the user to store a (ghost) object \verb`G`; each one allows the user to perform
a ghost operation called \emph{opening the invariant}, where they obtain temporary, exclusive
ownership over the \verb`G`. For example, this snippet shows how the implementation of
\verb`InvCell<T>::replace` temporarily gains access to the \verb`PermData<T>` object:

\begin{lstlisting}[language=Verus,style=VerusLineNos]
impl InvCell<T> {
    pub fn replace(&self, val: T) -> T {
        requires(self.wf() && self.inv(val));
        ensures(|old_val| self.inv(old_val));

        let r;

        // Opens the invariant `&self.perm_inv` which has type `LocalInvariant<PermData<T>>`.
        // Opening the invariant is a ghost operation, and it binds to the ghost variable `perm`.
        open_local_invariant!(&self.perm_inv => perm => {
            // The code inside, however, is executable. This is where we actually perform
            // the write, using ownership of the ghost `perm` object, of type `PermData<T>`.
            r = self.pcell.replace(&mut perm, val);
        });

        r // Return the old value.
    }
}
\end{lstlisting}

The difference between the two primitives is that \verb`LocalInvariant<G>` is restricted for use
on a single thread: it does not implement \verb`Send` or \verb`Sync`, the traits Rust uses to mark thread-safety.
\verb`AtomicInvariant<G>` is thread-safe, and it does implement these traits: however,
this comes with an additional requirement, that the invariant may only be opened
for atomic operations.
Since \verb`InvCell` (like Rust's \verb`Cell`) is for single-threaded use,
we use \verb`LocalInvariant<V>` here. 

The reader might wonder what happens if we attempt to nest calls to the invariant-opening
operation, \verb`open_local_invariant`. It would certainly be unsound if we could open the same
\verb`LocalInvariant<G>` object twice, and obtain double-ownership of the `T`.
Indeed, \verus generates extra verification conditions to disallow such things
by tracking which invariants are ``open'' at a given time. These verification conditions
are designed to be lightweight, and they have no impact on our SMT generation
for cases outside of those which use the low-level invariant APIs.


